%! Polynomial Long and Synthetic Division — Unit 3, Lesson 3.5 — AP Practice (ANSWER KEY)
% EXTRA CREDIT — optional, scored 1 to 5, exactly TWO pages, last in the packet (graded homework is DeltaMath).
% Page 1 is Section I, five multiple-choice items with four options (A)–(D); page 2 is Section
% II, one three-part free-response question on a table. Its contexts (a shipping crate, a
% startup's monthly profit) are used nowhere else in the lesson.
% Distractors are the lesson's errors on purpose: MC 1 uses c = +2 for x + 2 (A), skips the 0x
% placeholder (B), or takes the constant term as the remainder (D); MC 2 lets the remainder's
% degree reach or pass the divisor's (A, C, D); MC 3 uses c = +1 (A), keeps the dividend's degree
% (B), or copies coefficients without dividing (C); MC 4 flips the sign of a zero (B, C) or reads
% an output as a zero (D); MC 5 drops the placeholder (B), flips c (C), or both (A).
%
% PAGE LOCKSTEP with the other half of the pair: the key marks the correct option in its
% LABEL (no text added to the line); every \gradedblank line is identical in both files. The
% explicit \newpage fixes the page break in both files.
\documentclass[10pt]{article}
\usepackage{precalculus-article}
\usepackage{precalculus-key}
\usepackage{grading}
\gradingassignment{precalc-u03-l05-ap}
\gradingcourse{precalculus}
\newcommand{\TallMath}[1]{$\displaystyle #1\rule[-1.4em]{0pt}{3.2em}$}

\begin{document}

\pageheader{Unit 3, Lesson 3.5}{AP Practice --- Answer Key}
{\small\itshape Write each final answer in its box. Only what is inside a box is graded.}\par

\vspace{0.08in}

\boxguard[8]
\begin{remindbox}
\small
\textbf{Extra credit, optional, scored 1 to 5.} These two pages are written in the style of the AP
Precalculus exam. Your graded homework is on DeltaMath; do it first. Every
answer choice below is a mistake someone really makes, so read all four before you choose.
\end{remindbox}

\vspace{0.08in}

\begin{headlinebox}{goldbg}
\small
\textbf{Section I --- Multiple Choice.} Write the letter of the single best answer in its Letter box. No calculator.
\end{headlinebox}

\vspace{0.04in}

\begin{enumerate}[leftmargin=1.9em, itemsep=14pt]

\item What is the remainder when $p(x) = x^3 - 4x^2 + 6$ is divided by $x + 2$?

      \smallskip
      \textbf{(A)} $-2$ \hspace{2.6em} \textbf{(B)} $18$ \hspace{2.6em} {\color{keyred}\textbf{$\checkmark$\,(C)}} $-18$ \hspace{2.6em} \textbf{(D)} $6$ \hfill\textit{Letter:}~\gradedblank[id=ap1, type=choice, points=1, answer=C]{0.42in}{13pt}

\item A polynomial $f$ is divided by $g(x) = x^2 + 2$, so that $f(x) = (x^2 + 2)\,q(x) + r(x)$.
      Which of the following could be the remainder $r(x)$?

      \smallskip
      \textbf{(A)} $x^2 + 5$ \hspace{2.6em} {\color{keyred}\textbf{$\checkmark$\,(B)}} $3x - 4$ \hspace{2.6em} \textbf{(C)} $2x^2 - x$ \hspace{2.6em} \textbf{(D)} $x^3 + 1$ \hfill\textit{Letter:}~\gradedblank[id=ap2, type=choice, points=1, answer=B]{0.42in}{13pt}

\item A shipping crate has volume $V(x) = x^3 + 7x^2 + 14x + 8$ cubic feet, and its height is
      $x + 1$ feet. Which expression gives the area of the crate's base, in square feet?

      \smallskip
      \textbf{(A)} $x^2 + 8x + 22$ \hspace{1.6em} \textbf{(B)} $x^3 + 6x^2 + 8x$ \hspace{1.6em} \textbf{(C)} $x^2 + 7x + 14$ \hspace{1.6em} {\color{keyred}\textbf{$\checkmark$\,(D)}} $x^2 + 6x + 8$ \hfill\textit{Letter:}~\gradedblank[id=ap3, type=choice, points=1, answer=D]{0.42in}{13pt}

\item Selected values of a polynomial function $p$ are shown in the table. Which of the
      following must be a factor of $p(x)$?

      \smallskip
      {\renewcommand{\arraystretch}{1.25}%
      \begin{tabular}{|c|c|c|c|c|c|c|}
      \hline $x$ & $-2$ & $-1$ & $0$ & $1$ & $2$ & $3$ \\ \hline
      $p(x)$ & $-20$ & $0$ & $6$ & $4$ & $0$ & $0$ \\ \hline
      \end{tabular}}

      \smallskip
      {\color{keyred}\textbf{$\checkmark$\,(A)}} $x + 1$ \hspace{2.6em} \textbf{(B)} $x - 1$ \hspace{2.6em} \textbf{(C)} $x + 2$ \hspace{2.6em} \textbf{(D)} $x - 6$ \hfill\textit{Letter:}~\gradedblank[id=ap4, type=choice, points=1, answer=A]{0.42in}{13pt}

\item A student will use synthetic division to divide $2x^4 + 3x^2 - x + 5$ by $x + 2$. Which
      of the following is the correct setup?

      \smallskip
      {\setlength{\arraycolsep}{3pt}%
      \textbf{(A)} $\begin{array}{r|rrrr} 2 & 2 & 3 & -1 & 5 \end{array}$ \hspace{1.4em}
      \textbf{(B)} $\begin{array}{r|rrrr} -2 & 2 & 3 & -1 & 5 \end{array}$

      \smallskip
      \textbf{(C)} $\begin{array}{r|rrrrr} 2 & 2 & 0 & 3 & -1 & 5 \end{array}$ \hspace{1.4em}
      {\color{keyred}\textbf{$\checkmark$\,(D)}} $\begin{array}{r|rrrrr} -2 & 2 & 0 & 3 & -1 & 5 \end{array}$} \hfill\textit{Letter:}~\gradedblank[id=ap5, type=choice, points=1, answer=D]{0.42in}{13pt}

\end{enumerate}

\newpage

\begin{headlinebox}{goldbg}
\small
\textbf{Section II --- Free Response.} Show your reasoning outside the boxes. Answers about the
company must be written \emph{in context}, with units --- that is what earns credit on the AP
exam. Write each final answer in its box. Only what is inside a box is graded.
\end{headlinebox}

\vspace{0.04in}

\begin{enumerate}[leftmargin=1.9em, itemsep=6pt, start=6]

\item A new company's monthly profit is modeled by $P(x) = x^3 - 9x^2 + 23x - 15$, in thousands
      of dollars, $x$ months after launch, for $0 \le x \le 5$. Selected values are shown.

      \smallskip
      \begin{minipage}[t]{0.56\linewidth}
      \vspace{0pt}
      {\renewcommand{\arraystretch}{1.3}%
      \begin{tabular}{|c|c|c|c|c|c|c|}
      \hline $x$ (months) & $0$ & $1$ & $2$ & $3$ & $4$ & $5$ \\ \hline
      $P(x)$ (\$1000s) & $-15$ & $0$ & $3$ & $0$ & $-3$ & $0$ \\ \hline
      \end{tabular}}
      \end{minipage}\hfill
      \begin{minipage}[t]{0.4\linewidth}
      \vspace{0pt}
      {\footnotesize\itshape Work for (a):}
      \vspace{0.55in}
      \end{minipage}

      \textbf{(a)} The table shows $P(1) = 0$. Explain why $x - 1$ must be a factor of $P(x)$, and
      use synthetic division to write $P(x) = (x - 1)\cdot q(x)$.
      \rubricbox[2pt]{\gradedblank[id=ap6a, type=rubric, points=2,
        display={Because $P(1) = 0$, the Factor Theorem says $x - 1$ is a factor (dividing by $x - 1$ leaves remainder $P(1) = 0$). Synthetic division with $c = 1$ on $1,\ -9,\ 23,\ -15$ gives $1,\ -8,\ 15 \mid 0$, so $P(x) = (x - 1)(x^2 - 8x + 15)$.}]{\linewidth}{2in}}

      \textbf{(b)} Factor $q(x)$ and find every time the company breaks even (profit of \$0).
      Interpret your answer in context.
      \rubricbox[2pt]{\gradedblank[id=ap6b, type=rubric, points=2,
        display={$q(x) = x^2 - 8x + 15 = (x - 3)(x - 5)$, so $P(x) = (x - 1)(x - 3)(x - 5)$. The company breaks even (profit \$0) at 1, 3, and 5 months after launch.}]{\linewidth}{2in}}

      \textbf{(c)} An analyst divides $P(x)$ by $x - 4$ and reports a remainder of $3$. Use the table
      to decide whether the analyst is right, and interpret the correct remainder in context.
      \rubricbox[2pt]{\gradedblank[id=ap6c, type=rubric, points=2,
        display={The analyst is wrong. By the Remainder Theorem the remainder is $P(4) = -3$, read from the table. Four months after launch the company's profit is $-3$ thousand dollars: it loses \$3000 that month. (3 is $P(2)$, not $P(4)$.)}]{\linewidth}{2in}}

\end{enumerate}

\end{document}
